\documentclass{article}
\usepackage[utf8]{inputenc}
\usepackage{tabularx}
\usepackage{booktabs}
\usepackage[margin=1in]{geometry}
\usepackage{enumitem}
\usepackage{array}
\usepackage{lmodern}

\newcommand{\usecasetable}[9]{
	\begin{table}[h!]
		\centering
		\footnotesize
		\renewcommand{\arraystretch}{1.4}
		\begin{tabularx}{\textwidth}{|l|>{\raggedright\arraybackslash}X|}
			\hline
			\textbf{ID Use Case} & #1 \\ \hline
			\textbf{Nome Use Case} & #2 \\ \hline
			\textbf{Attori} & #3 \\ \hline
			\textbf{Descrizione} & #4 \\ \hline
			\textbf{Precondizioni} & #5 \\ \hline
			\textbf{Postcondizioni} & #6 \\ \hline
			\textbf{Flusso Principale} & #7 \\ \hline
			\textbf{Flussi alternativi} & #8 \\ \hline
			\textbf{Business Rules e Riferimenti} & #9 \\ \hline
		\end{tabularx}
	\end{table}
}

\begin{document}
	\usecasetable
	{UC-01}
	{Login to Account}
	{User (Admin, Customer, Owner)}
	{Permette ad un utente registrato di accedere all'account. Gestisce anche il cambio password obbligatorio per i nuovi Owner e Admin.}
	{L'utente possiede credenziali valide registrate nel sistema.}
	{L'utente viene autenticato e reindirizzato alla propria dashboard (in base al ruolo).}
	{
		\begin{enumerate}[nosep, leftmargin=*]
			\item L'utente inserisce Email (o Username) e Password nella pagina di login;
			\item L'utente clicca su "Accedi";
			\item Il sistema verifica che le credenziali siano corrette;
			\item Il sistema verifica lo stato dell'account (deve essere ATTIVO e non bannato);
			\item Il sistema controlla se l'account richiede un cambio password obbligatorio (flag OTP);
			\item Il sistema concede l'accesso e crea la sessione utente.
		\end{enumerate}
	}
	{
		\begin{enumerate}[nosep, leftmargin=25px]
			\item[\textbf{3a-e.}] \textbf{Credenziali Errate:} Il sistema mostra un messaggio di errore ("Email o password non valide") e invita a riprovare.
			\item[\textbf{4a-e.}] \textbf{Utente Bannato - App Ban:} Il sistema rileva che l'utente ha subito un "App Ban" dall'Admin. Il sistema blocca l'accesso e mostra il messaggio: "Il tuo account è stato disabilitato per violazione dei termini.".
			\item[\textbf{4b-e.}] \textbf{Account Customer chiuso:} Il sistema rileva che l'account è stato chiuso; blocca quindi l'accesso e mostra il messaggio:"Impossibile accedere: account chiuso.".
			\item[\textbf{4c-e.}] \textbf{Spiaggia Disattivata - Solo Owner:} Se l'utente è un Owner e la sua spiaggia è stata disattivata (Ban Spiaggia), il sistema blocca l'accesso e mostra: "La tua struttura è stata sospesa dall'amministratore.".
			\item[\textbf{5a.}] \textbf{Cambio Password Obbligatorio (Primo Accesso Owner o Admin):} Se il sistema rileva il flag temporaneo sulla password, sospende il login e reindirizza l'utente a una schermata obbligatoria di "Imposta Nuova Password". 
			\begin{enumerate}[nosep, leftmargin=*]
				\item[i.] L'utente inserisce la nuova password e la conferma;
				\item[ii.] Il sistema la salva crittografata, rimuove il flag temporaneo e procede col normale login (torna al passo 6).
			\end{enumerate}
		\end{enumerate}
	}
	{
		\begin{enumerate}[nosep, leftmargin=*]
			\item[•] Un account con "App Ban" non può mai generare una sessione di login valida.
			\item[•] I nuovi account Owner e Admin creati da un (altro) Admin non possono accedere alle funzionalità del sistema senza aver prima sostituito la password OTP.
		\end{enumerate}
		
		\begin{enumerate}[leftmargin=*]
			\item[•] \textbf{Testing}: le classi che si occupano dei test sono \texttt{AuthenticationServiceTest} nel package \texttt{service}; \texttt{UserTest}, \texttt{CustomerTest} \texttt{OwnerTest} e \texttt{AdminTest} nel package \texttt{domain}.
		\end{enumerate}
	}
	
	\newpage
	
	\usecasetable
	{UC-02}
	{Manage Account}
	{User (Admin, Customer, Owner)}
	{Permette all'utente di modificare i suoi dati legati all'account (username, email, numero di telefono, indirizzo ecc.).}
	{L'utente in questione deve aver fatto login inizialmente all'account.}
	{I dati modificati vengono aggiornati nel database e viene mostrato all'utente un messaggio di conferma dei cambiamenti.}
	{
		\begin{enumerate}[nosep, leftmargin=*]
			\item[1.] L'utente naviga dentro la pagina "Impostazioni Account/Modifica infornazioni";
			\item[2.] Il sistema recupera i dati e mostra i dati attuali dell'account (nome, cognome);
			\item[3.] L'utente modifica i dati che preferisce;
			\item[4.] L'utente clicca su "Salva modifiche";
			\item[5.] Il sistema convalida e salva i nuovi dati;
			\item[6.] Viene mosrato all'utente un messaggio di conferma.
		\end{enumerate}
	}
	{
		\begin{enumerate}[nosep, leftmargin=25px]
			\item[\textbf{2a.}] \textbf{L'utente è un Customer:} Il sistema mostra anche l'indirizzo;
			\item[\textbf{3a.}] \textbf{L'utente è un Customer:} L'utente può modificare anche l'indirizzo; 
			\item[\textbf{1a.}] \textbf{Cambio email (da passo 1):}
			\begin{enumerate}[nosep, leftmargin=*]
				\item[i.] L'utente naviga su "Cambia Email"; 
				\item[ii.] Il sistema mostra una pagina dedicata per il cambio email;
				\item[iii.] L'utente inserisce la nuova email e la password attuale; 
				\item[iv.] Il sistema verifica l'unicità della email nel database e aggiorna la nuova email nel database.
			\end{enumerate}
			\item[\textbf{1b-e.}] \textbf{Email già usata per un altro account/non valida:} Il sistema segnala l'errore all'utente e invita a ricontrollare.
			\item[\textbf{1b-e.}] \textbf{Password non valida:} Il sistema segnala l'errore all'utente e invita a reinserire la password.
			\item[\textbf{1c.}] \textbf{Cambio numero di telefono (L'utente è un Customer, da passo 1):} 
			\begin{enumerate}[nosep, leftmargin=*]
				\item[i.] L'utente naviga su "Cambia numero di telefono";
				\item[ii.] Il sistema mostra una pagina dedicata per il cambio numero di telefono;
				\item[iii.] L'utente inserisce il nuovo numero di telefono e la password attuale;
				\item[iv.] Il sistema verifica l'unicità del numero nel database e aggiorna il nuovo numero nel database.
			\end{enumerate}
			\item[\textbf{1c-e.}] \textbf{Numero di telefono già usato per un altro account/non valido:} Il sistema segnala l'errore all'utente e invita a ricontrollare.
			\item[\textbf{1c-e.}] \textbf{Password non valida:} Il sistema segnala l'errore all'utente e invita a reinserire la password.
			\item[\textbf{1d.}] \textbf{Cambio password (da passo 1):}
			\begin{enumerate}[nosep, leftmargin=1px]
				\item[i.] L'utente naviga su "Cambia password"; 
				\item[ii.] Il sistema mostra una pagina dedicata per il cambio password;
				\item[iii.] L'utente inserisce la nuova e la vecchia password;
				\item[iv.] Il sistema verifica la password vecchia;
				\item[v.] Il sistema aggiorna la password nel database e mostra un messaggio di conferma all'utente.
			\end{enumerate}
			\item[\textbf{1d-e.}] \textbf{Password vecchia non valida:} Il sistema segnala l'errore all'utente e invita a reinserire la vecchia password.
		\end{enumerate}
	}
	{
		\begin{enumerate}[nosep, leftmargin=*]
			\item[•] Ogni utente deve inserire dati validi. I dati non possono essere nulli.
			\item[•] Solo Customer possiede indirizzo e numero di telefono;
			\item[•] Username, nome, cognome (per i Customers: indirizzo) possono essere cambiati allo stesso momento. Email, Password (per i Customers: numero di telefono) invece devono essere cambiati singolarmente attraversando dei check di sicurezza.
		\end{enumerate}
		
		
		\begin{enumerate}[leftmargin=*]
			\item[•] \textbf{Testing}: le classi che si occupano dei test sono \texttt{ManageUserServiceTest}, \texttt{ManageCustomerServiceTest}, \texttt{ManageOwnerServiceTest} e \texttt{ManageAdminServiceTest} nel package \texttt{service}; \texttt{UserTest}, \texttt{CustomerTest}, \texttt{OwnerTest} e \texttt{AdminTest} nel package \texttt{domain}. \newline Vengono usate anche \texttt{service.AddressServiceTest} e \texttt{domain.AddressTest} in caso di un \texttt{Customer} che decide di modificare l'indirizzo.
		\end{enumerate}
	}
	
	\newpage
	
	\usecasetable
	{UC-03}
	{Create Customer Account}
	{Guest}
	{Permette a un utente non registrato di creare un nuovo profilo Customer inserendo i dati anagrafici e verificando i propri recapiti (Email e Telefono).}
	{L'utente non deve essere autenticato nel sistema (stato Guest).}
	{Viene creato un nuovo record nel database per l'account Customer.}
	{
		\begin{enumerate}[nosep, leftmargin=*]
			\item L'utente accede alla pagina di registrazione;
			\item Il sistema mostra il modulo di inserimento dati (Username, Nome, Cognome, Indirizzo, Numero di Telefono, Email, Password);
			\item L'utente inserisce tutti i dati richiesti e preme "Registrati";
			\item Il sistema verifica che lo Username non sia già presente nel database;
			\item Il sistema verifica la validità di Email e Numero di Telefono;
			\item Il sistema attiva l'account e mostra un messaggio di successo.
		\end{enumerate}
	}
	{
		\begin{enumerate}[nosep, leftmargin=25px]
			\item[\textbf{4a.}] \textbf{Username già esistente:} Il sistema segnala che lo username è già utilizzato e invita a sceglierne uno nuovo.
			\item[\textbf{5a.}] \textbf{Numero di telefono/Email già esistente:} Se l'email o il telefono sono già associati ad un account, il sistema evidenzia i campi errati.
		\end{enumerate}
	}
	{
		\begin{enumerate}[nosep, leftmargin=*]
			\item[•] Tutti i campi sono obbligatori (non null).
			\item[•] Lo Username deve essere univoco a livello globale.
		\end{enumerate}
		
		\begin{enumerate}[leftmargin=*]
			\item[•] \textbf{Testing}: le classi che si occupano dei test sono \texttt{CreateUserServiceTest}, \texttt{CreateCustomerServiceTest} e \texttt{AddressServiceTest} nel package \texttt{service}; \texttt{UserTest}, \texttt{CustomerTest} e \texttt{AddressTest} nel package \texttt{domain}.
		\end{enumerate}
	}
		
	\newpage
	
	\usecasetable
	{UC-04}
	{Browse Active Beaches}
	{Guest, Customer}
	{Permette agli utenti di visualizzare l'elenco delle spiagge operative, effettuare ricerche mirate per nome o città e consultare i dettagli di ogni stabilimento.}
	{Nessuna (Accessibile sia ad utenti Guest che Customer). Per interagire come Customer (es. vedere il tasto prenota), l'utente deve aver effettuato il login e il suo account deve essere in stato ATTIVO (non chiuso, non disattivato e non soggetto a Ban).}
	{L'utente ha visualizzato le informazioni di una spiaggia; il Customer, oltre a ciò, ha anche la possibilità di avviare una prenotazione.}
	{
		\begin{enumerate}[nosep, leftmargin=*]
			\item L'utente accede alla sezione "Esplora Spiagge" dell'applicazione;
			\item Il sistema recupera dal database e mostra l'elenco completo di tutte le spiagge con stato "Attivo";
			\item L'utente seleziona una spiaggia dalla lista;
			\item Il sistema mostra la pagina di dettaglio della spiaggia (Descrizione, Servizi, Posizione, Valutazione media e Lista delle Recensioni degli utenti).
		\end{enumerate}
	}
	{
		\begin{enumerate}[nosep, leftmargin=25px]
			\item[\textbf{2a.}] \textbf{Ricerca per Nome o Città:} 
			\begin{enumerate}[nosep, leftmargin=*]
				\item[i.] L'utente inserisce il nome dello stabilimento o il nome della città nella barra di ricerca;
				\item[ii.] Il sistema filtra i risultati e mostra solo le spiagge corrispondenti ai criteri;
				\item[iii.] Se non ci sono risultati, il sistema mostra il messaggio "Nessuna spiaggia trovata".
			\end{enumerate}
			\item[\textbf{4a.}] \textbf{Interazione Customer (Prenotazione):}
			\begin{enumerate}[nosep, leftmargin=*]
				\item[i.] Se l'utente è autenticato come Customer, il sistema mostra il pulsante "Prenota!";
				\item[ii.] Cliccando, il sistema reindirizza l'utente a UC-05 (Create Booking).
			\end{enumerate}
			\item[\textbf{4b.}] \textbf{Interazione Guest:} 
			\begin{enumerate}[nosep, leftmargin=*]
				\item[i.] Se l'utente è Guest, il sistema mostra, oltre alle informazioni della spiaggia, un messaggio informativo: "Accedi o Registrati per poter effettuare una prenotazione".
			\end{enumerate}
		\end{enumerate}
	}
	{
		\begin{enumerate}[nosep, leftmargin=*]
			\item[•] Vengono mostrate esclusivamente le spiagge il cui stato è "Attivo" (impostato dall'Owner).
			\item[•] I risultati della ricerca per città devono includere anche spiagge limitrofe se configurato nel sistema.
			\item[•] La visualizzazione dei dettagli non impegna alcuna risorsa (non blocca posti spiaggia).
		\end{enumerate}
		
		\begin{enumerate}[leftmargin=*]
			\item[•] \textbf{Mockup}: si vedano Figure 16 e 17 - Esplora Spiagge (Desktop e Mobile)
			
			\item[•] \textbf{Testing}: le classi che si occupano dei test sono \texttt{BrowseBeachServiceTest} e \texttt{ViewBeachReviewsServiceTest} nel package \texttt{service}; \texttt{BeachTest}, \texttt{BeachGeneralTest}, \texttt{BeachServicesTest} e \texttt{ReviewTest} nel package \texttt{domain}.
		\end{enumerate}
	}
		
	\newpage
	
	\usecasetable
	{UC-05}
	{Create Booking}
	{Customer, Owner}
	{Permette ad un Customer di prenotare posti, servizi extra e parcheggi presso una spiaggia in una certa data. Permette inoltre ad un Owner di registrare manualmente una prenotazione per un cliente telefonico.}
	{L'utente deve aver effettuato il login. Se l'utente è un Customer, non deve essere bannato dalla spiaggia selezionata. Se l'utente è un Owner, deve possedere una spiaggia attiva. La spiaggia selezionata (o posseduta dall'Owner) deve essere nello stato ATTIVO.}
	{Viene creato un nuovo record di prenotazione nel database con un prezzo totale calcolato e uno stato iniziale (PENDING per i Customer, CONFIRMED per gli Owner).}
	{
		\begin{enumerate}[nosep, leftmargin=*]
			\item Il Customer seleziona una spiaggia (partendo da UC-04) e sceglie una data;
			\item Il sistema verifica che il Customer non sia bannato da quella spiaggia;
			\item Il sistema mostra la mappa o la lista delle disponibilità per la data scelta;
			\item Il Customer seleziona i posti spiaggia desiderati (es. ombrelloni, lettini);
			\item Il Customer seleziona eventuali extra (lettini, sdraio, sedie);
			\item Il Customer inserisce il numero di parcheggi desiderati;
			\item Il sistema calcola e mostra il prezzo totale riepilogativo;
			\item Il Customer conferma la prenotazione;
			\item Il sistema salva la prenotazione e le assegna lo stato "PENDING", mostrando un messaggio di successo.
		\end{enumerate}
	}
	{
		\begin{enumerate}[nosep, leftmargin=36px]
			\item[\textbf{1a.}] \textbf{Flusso Owner (Prenotazione Telefonica):}
			\begin{enumerate}[nosep, leftmargin=*]
				\item[i.] L'Owner accede al pannello gestionale della propria spiaggia e avvia una nuova prenotazione;
				\item[ii.] L'Owner seleziona la data, i posti, gli extra e i parcheggi (come nei passi 3-6 del Main Flow);
				\item[iii.] L'Owner inserisce obbligatoriamente il Nome e il Numero di Telefono del cliente chiamante;
				\item[iv.] Il sistema calcola e mostra il totale;
				\item[v.] L'Owner conferma la prenotazione;
				\item[vi.] Il sistema salva la prenotazione assegnandole automaticamente lo stato "CONFIRMED".
			\end{enumerate}
			\item[\textbf{2a-e.}] \textbf{Utente Bannato:} Il sistema rileva che il Customer è nella blacklist della spiaggia e blocca l'operazione mostrando il messaggio: "Non sei autorizzato a prenotare in questa struttura.".
			\item[\textbf{4/6a-e.}] \textbf{Disponibilità Esaurita:} Se durante la selezione i posti o i parcheggi scelti vengono esauriti (es. prenotati da un altro utente nello stesso istante), il sistema avvisa l'utente e lo invita a modificare la scelta.
		\end{enumerate}
	}
	{
		\begin{enumerate}[nosep, leftmargin=*]
			\item[•] Le prenotazioni dei Customer partono sempre dallo stato PENDING e richiedono un passaggio a CONFIRMED da parte dell'Owner.
			\item[•] Le prenotazioni create dall'Owner sono considerate già confermate (CONFIRMED). 
			\item[•] L'Owner non può selezionare la spiaggia: il sistema forza la prenotazione solo sullo stabilimento di sua proprietà.
			\item[•] Il calcolo del totale si basa sul tariffario (Season/Zone) impostato per quella data.
		\end{enumerate}
		
		\begin{enumerate}[leftmargin=*]
			\item[•] \textbf{Mockup}: si veda Figura 16 - Flusso di Prenotazione
			
			\item[•] \textbf{Testing}: le classi che si occupano dei test sono \texttt{CreateBookingServiceTest} e \texttt{CreateManualBookingServiceTest} nel package \texttt{service}; \texttt{BookingTest}, \texttt{BookingParkingTest} e \texttt{PriceCalculatorTest} nel package \texttt{domain}.
		\end{enumerate}
	}
		
	\newpage
	
	\usecasetable
	{UC-06}
	{Manage Bookings}
	{Customer, Owner}
	{Permette ai Customer di modificare o cancellare le proprie prenotazioni e agli Owner di approvare/rifiutare le richieste in sospeso o gestire le prenotazioni telefoniche.}
	{L'utente deve aver effettuato il login. Deve esistere almeno una prenotazione associata all'utente (Customer) o alla spiaggia (Owner).}
	{La prenotazione viene aggiornata e il suo stato (PENDING, CONFIRMED, REJECTED, CANCELLED) viene modificato di conseguenza nel database.}
	{
		\begin{enumerate}[nosep, leftmargin=*]
			\item L'utente accede alla sezione "Le mie prenotazioni" (Customer) o "Gestione Prenotazioni" (Owner);
			\item Il sistema recupera e mostra lo storico delle prenotazioni;
			\item L'utente seleziona una prenotazione specifica per visualizzarne i dettagli;
			\item Il sistema verifica i vincoli temporali (la data odierna deve essere precedente alla data della prenotazione) e lo stato attuale;
			\item L'utente (Customer) sceglie di modificare i parametri (aggiungere/rimuovere Posti, Extra, Parcheggi) oppure di "Cancellare" la prenotazione;
			\item Nel caso di modifica, il sistema ricalcola il totale e l'utente conferma;
			\item Il sistema salva l'aggiornamento: se la prenotazione modificata era CONFIRMED, il sistema la declassa allo stato PENDING. Se l'utente ha scelto di cancellare, lo stato diventa CANCELLED.
		\end{enumerate}
	}
	{
		\begin{enumerate}[nosep, leftmargin=25px]
			\item[\textbf{5a.}] \textbf{Flusso Owner (Valutazione prenotazioni stato PENDING):} 
			\begin{enumerate}[nosep, leftmargin=*]
				\item[i.] L'Owner seleziona una prenotazione Customer in stato PENDING;
				\item[ii.] L'Owner clicca su "Approva" o "Rifiuta";
				\item[iii.] Il sistema aggiorna lo stato in CONFIRMED o REJECTED e notifica il Customer (tramite mail e/o SMS).
			\end{enumerate}
			\item[\textbf{5b.}] \textbf{Flusso Owner (Modifica Prenotazione Telefonica):} 
			\begin{enumerate}[nosep, leftmargin=*]
				\item[i.] L'Owner seleziona una prenotazione registrata manualmente;
				\item[ii.] Applica le modifiche (posti, extra, parcheggi) o cancella, rispettando gli stessi vincoli di tempo del Customer.
			\end{enumerate}
			\item[\textbf{4a-e.}] \textbf{Vincolo Temporale Scaduto:} Se la data odierna è uguale o successiva alla data della prenotazione, il sistema nasconde/disabilita i pulsanti di modifica e cancellazione.
			\item[\textbf{4b-e.}] \textbf{Stato Non Modificabile:} Se la prenotazione è in stato REJECTED o CANCELLED, il sistema impedisce qualsiasi modifica da parte del Customer/Owner.
		\end{enumerate}
	}
	{
		\begin{enumerate}[nosep, leftmargin=*]
			\item[•] La Spiaggia e la Data della prenotazione sono immutabili: per cambiarle, l'utente deve cancellare la prenotazione e crearne una nuova.
			\item[•] Ogni modifica ai posti o ai servizi da parte del Customer riporta una prenotazione CONFIRMED allo stato PENDING (richiedendo una nuova approvazione da parte dell'Owner).
			\item[•] Le modifiche sono inibite a partire dalla mezzanotte ($00:00$) del giorno stesso della prenotazione.
		\end{enumerate}
		
		\begin{enumerate}[leftmargin=*]
			\item[•] \textbf{Mockup}: si vedano Figure 17 e 18 - Gestione Prenotazioni (Customer e Owner)
			
			\item[•] \textbf{Testing}: le classi che si occupano dei test sono \texttt{BookingServiceTest} nel package \texttt{service}; \texttt{BookingTest}, \texttt{BookingParkingTest} e \texttt{PriceCalculatorTest} nel package \texttt{domain}.
		\end{enumerate}
	}
		
	\newpage
	
	\usecasetable
	{UC-07}
	{Manage Reviews}
	{Customer}
	{Permette ad un Customer di valutare e commentare la propria esperienza presso una spiaggia in cui ha effettivamente soggiornato, oppure di eliminare una propria recensione precedentemente pubblicata.}
	{L'utente è loggato. La spiaggia è in stato ATTIVO. L'utente non ha un Ban per quella spiaggia. Deve esistere almeno una prenotazione per quella spiaggia con data nel passato (rispetto a quando viene applicato lo Use Case) e in stato CONFIRMED.}
	{La recensione viene salvata o rimossa dal database, associata alla spiaggia e all'utente.}
	{
		\begin{enumerate}[nosep, leftmargin=*]
			\item L'utente naviga nella sezione "Le mie prenotazioni" o sulla pagina della spiaggia;
			\item L'utente clicca sul pulsante "Lascia una recensione" per la spiaggia selezionata;
			\item Il sistema verifica che tutti i requisiti siano soddisfatti (prenotazione passata e CONFIRMED, assenza di ban, spiaggia attiva);
			\item Il sistema mostra il modulo della recensione;
			\item L'utente seleziona una valutazione (da 1 a 5 stelle) e inserisce un commento testuale;
			\item L'utente clicca su "Pubblica";
			\item Il sistema salva la recensione nel database e mostra un messaggio di conferma.
		\end{enumerate}
	}
	{
		\begin{enumerate}[nosep, leftmargin=25px]
			\item[\textbf{2a.}] \textbf{Eliminazione Recensione:} 
			\begin{enumerate}[nosep, leftmargin=*]
				\item[i.] L'utente seleziona una recensione che ha già pubblicato in passato e clicca su "Elimina";
				\item[ii.] Il sistema chiede conferma dell'operazione ("Sei sicuro di voler eliminare questa recensione?");
				\item[iii.] L'utente conferma;
				\item[iv.] Il sistema rimuove il record dal database e mostra un messaggio di successo.
			\end{enumerate}
			\item[\textbf{3a-e.}] \textbf{Mancanza Requisiti:} Se l'utente tenta di accedere all'URL della recensione senza averne diritto (es. prenotazione futura, prenotazione CANCELLED/REJECTED o mai stato in quella spiaggia), il sistema blocca l'azione mostrando: "Devi aver completato un soggiorno in questa struttura per poter lasciare una recensione.".
			\item[\textbf{3b-e.}] \textbf{Spiaggia Disattivata/Ban:} Se la spiaggia non è più attiva o l'utente è stato bannato da quella struttura dopo il soggiorno, il sistema inibisce la funzione di recensione.
			\item[\textbf{6a-e.}] \textbf{Recensione non completata:} Se l'utente non compila entrambi i campi, il sistema evidenzia i campi da compilare obbligatoriamente.
		\end{enumerate}
	}
	{
		\begin{enumerate}[nosep, leftmargin=*]
			\item[•] La valutazione è obbligatoria ed è un valore intero tra 1 e 5.
			\item[•] Il commento testuale è anch'esso obbligatorio.
			\item[•] Il sistema accetta recensioni solo per prenotazioni la cui data è strettamente minore della data odierna (es. prenotazione il 19 luglio, recensibile dal 20 luglio).
			\item[•] L'App Ban è già gestito a monte dal Login (UC-01).
		\end{enumerate}
		
		\begin{enumerate}[leftmargin=*]
			\item[•] \textbf{Mockup}: si veda Figura 17 - Le Mie Prenotazioni e Recensioni
			
			\item[•] \textbf{Testing}: le classi che si occupano dei test sono \texttt{ReviewServiceTest} nel package \texttt{service}; \texttt{ReviewTest} nel package \texttt{domain}.
		\end{enumerate}
	}
		
	\newpage
	
	\usecasetable
	{UC-08}
	{Close Account}
	{Customer}
	{Permette ad un Customer di chiudere definitivamente il proprio profilo. L'operazione comporta la cancellazione automatica di tutte le prenotazioni future.}
	{L'utente deve aver effettuato il login.}
	{L'account del Customer viene disattivato permanentemente. Tutte le sue prenotazioni future passano allo stato CANCELLED. La sessione corrente viene terminata.}
	{
		\begin{enumerate}[nosep, leftmargin=*]
			\item Il Customer accede alla sezione "Impostazioni Account" e seleziona l'opzione "Chiudi Account";
			\item Il sistema mostra un avviso critico (Warning) che spiega le conseguenze: perdita definitiva dell'accesso e cancellazione di tutte le prenotazioni future;
			\item Il sistema richiede l'inserimento della password attuale per confermare l'identità;
			\item Il Customer inserisce la password e conferma l'intenzione di eliminare l'account;
			\item Il sistema valida la password;
			\item Il sistema individua tutte le prenotazioni del Customer (in stato PENDING o CONFIRMED) con data successiva a quella odierna;
			\item Il sistema imposta lo stato di queste prenotazioni su CANCELLED;
			\item Il sistema disattiva definitivamente l'account del Customer;
			\item Il sistema esegue il logout forzato e reindirizza l'utente alla homepage.
		\end{enumerate}
	}
	{
		\begin{enumerate}[nosep, leftmargin=25px]
			\item[\textbf{4a.}] \textbf{Annullamento operazione:} Se il Customer clicca su "Annulla" o chiude la finestra di avviso, il processo si interrompe e il sistema lo riporta alle impostazioni.
			\item[\textbf{5a-e.}] \textbf{Password errata:} Se la password inserita per la conferma non è corretta, il sistema blocca l'eliminazione, mostra un messaggio di errore e invita a riprovare.
		\end{enumerate}
	}
	{
		\begin{enumerate}[nosep, leftmargin=*]
			\item[•] Le prenotazioni passate (storico) non vengono cancellate, ma restano a disposizione degli Owner per motivi fiscali e di reportistica.
			\item[•] Vengono cancellate solo le prenotazioni con data $>$ data odierna.
			\item[•] L'operazione è irreversibile: il Customer non potrà più accedere né recuperare l'account.
		\end{enumerate}
		
		\begin{enumerate}[leftmargin=*]
			\item[•] \textbf{Testing}: le classi che si occupano dei test sono \texttt{ManageCustomerServiceTest} nel package \texttt{service}; \texttt{UserTest} e \texttt{CustomerTest} nel package \texttt{domain}.
		\end{enumerate}
	}
		
	\newpage
	
	\usecasetable
	{UC-09}
	{Create Report}
	{Customer, Owner}
	{Permette ad un Customer di segnalare una spiaggia, oppure ad un Owner di segnalare un cliente, a seguito di un'interazione reale e conclusa.}
	{L'utente deve aver effettuato il login. La spiaggia deve essere in stato ATTIVO. Deve esistere una prenotazione in stato CONFIRMED con data strettamente precedente a quella odierna. Il Customer non deve essere attualmente bannato (né dall'App né dalla spiaggia) e non deve avere l'account chiuso.}
	{La segnalazione viene salvata nel database con stato PENDING, in attesa della revisione da parte di un Admin.}
	{
		\begin{enumerate}[nosep, leftmargin=*]
			\item Il Customer accede alla sezione delle proprie prenotazioni;
			\item Seleziona una prenotazione passata relativa a una specifica spiaggia;
			\item Clicca sul pulsante "Crea Segnalazione";
			\item Il sistema verifica i requisiti (spiaggia attiva, prenotazione CONFIRMED, data passata, utente non bannato);
			\item Il Customer digita il testo descrittivo del problema riscontrato;
			\item Il Customer clicca su "Invia";
			\item Il sistema salva il report impostando lo stato su PENDING e mostra un messaggio di conferma e mandata una email al Customer con il Report creato.
		\end{enumerate}
	}
	{
		\begin{enumerate}[nosep, leftmargin=25px]
			\item[\textbf{1a.}] \textbf{Flusso Owner (Segnala Customer):} 
			\begin{enumerate}[nosep, leftmargin=*]
				\item[i.] L'Owner accede alla sezione "Crea Segnalazione" del proprio pannello gestionale;
				\item[ii.] Il sistema mostra la lista di tutte le prenotazioni passate in stato CONFIRMED relative alla sua spiaggia;
				\item[iii.] L'Owner seleziona una prenotazione/cliente dalla lista e clicca su "Segnala";
				\item[iv.] Il sistema esegue i controlli di validità (incluso verificare che il Customer non sia già stato bannato);
				\item[v.] L'Owner digita il testo e clicca su "Invia";
				\item[vi.] Il sistema salva il report in stato PENDING e manda una copia del Report alla email dell'Owner.
			\end{enumerate}
			\item[\textbf{4a-e.}] \textbf{Requisiti non soddisfatti:} Se le condizioni non sono rispettate (es. prenotazione futura, stato diverso da CONFIRMED, spiaggia inattiva, o Customer già bannato), il sistema disabilita la funzione di segnalazione o mostra un messaggio di errore: "Azione non consentita per questa prenotazione".
			\item[\textbf{6a-e.}] \textbf{Dati non compilati:} Se l'utente non compila il campo descrizione del problema, il sistema evidenzia il campo invitando l'utente a compilarlo.
		\end{enumerate}
	}
	{
		\begin{enumerate}[nosep, leftmargin=*]
			\item[•] Ogni segnalazione viene inizializzata in stato PENDING ed è visibile unicamente agli Admin.
			\item[•] Non è possibile segnalare un Customer se le misure restrittive (App Ban o Beach Ban) sono già attive su di lui.
			\item[•] La regola temporale è rigorosa: si possono segnalare solo eventi per cui la Data Prenotazione $<$ Data Odierna.
		\end{enumerate}
		
		\begin{enumerate}[leftmargin=*]
			\item[•] \textbf{Testing}: le classi che si occupano dei test sono \texttt{CreateReportServiceTest} nel package \texttt{service}; \texttt{ReportTest} nel package \texttt{domain}.
		\end{enumerate}
	}
		
	\newpage
	
	\usecasetable
	{UC-10}
	{Create Beach}
	{Owner}
	{Permette a un Owner appena registrato (o senza strutture associate) di inserire nel sistema la propria spiaggia, configurandone dati anagrafici, inventario, servizi, parcheggi, stagioni e zone.}
	{L'utente deve aver effettuato il login come Owner. Nel database NON deve esistere alcuna spiaggia già associata a questo account.}
	{Viene creato un nuovo record "Spiaggia" nel database, associato all'Owner. Lo stato sarà ATTIVO se configurata completamente e confermata, altrimenti INATTIVO (bozza).}
	{
		\begin{enumerate}[nosep, leftmargin=15px]
			\item L'Owner accede al sistema; il sistema rileva l'assenza di strutture e mostra il pulsante "Crea Spiaggia";
			\item L'Owner clicca sul pulsante e accede al modulo di creazione;
			\item L'Owner inserisce i Dati Obbligatori: Nome, Descrizione, Numero di Telefono, Indirizzo e Informazioni Extra;
			\item L'Owner (opzionalmente) compila l'Inventario: quantità di ombrelloni, tende, sdraio, lettini, sedie, camerini;
			\item L'Owner (opzionalmente) seleziona i Servizi: bagni, docce, piscina, bar, ristorante, Wi-Fi, campo da pallavolo;
			\item L'Owner (opzionalmente) configura il Parcheggio: numero posti auto, moto, bici, veicoli elettrici e presenza telecamere di videosorveglianza;
			\item L'Owner (opzionalmente) crea le Stagioni (prezzi base e tariffe sulle varie zone) e le Zone (layout ombrelloni/tende);
			\item L'Owner clicca su "Salva Spiaggia";
			\item Il sistema convalida i dati obbligatori;
			\item Il sistema verifica la completezza della configurazione (presenza sia dei dati obbligatori che di quelli opzionali/operativi);
			\item Avendo inserito tutti i dati, il sistema mostra un prompt: "Vuoi attivare la spiaggia e renderla visibile ora?";
			\item L'Owner accetta (o rifiuta);
			\item Il sistema salva definitivamente la spiaggia nel database (in stato ATTIVO se l'Owner ha accettato, altrimenti INATTIVO).
		\end{enumerate}
	}
	{
		\begin{enumerate}[nosep, leftmargin=25px]
			\item[\textbf{9a-e.}] \textbf{Campi Obbligatori Mancanti:} Se Nome, Indirizzo o altri campi base sono vuoti, il sistema evidenzia gli errori e impedisce il salvataggio.
			\item[\textbf{10a.}] \textbf{Salvataggio Parziale (Bozza):} Se l'Owner ha inserito i dati obbligatori ma ha saltato configurazioni chiave (es. Stagioni e Zone), il sistema non mostra il prompt di attivazione. Salva direttamente la spiaggia in stato INATTIVO. L'Owner dovrà usare "Manage Beach" (UC-11) in futuro per completarla e attivarla.
		\end{enumerate}
	}
	{
		\begin{enumerate}[nosep, leftmargin=*]
			\item[•] Relazione 1:1 rigida: un Owner può possedere e gestire al massimo una spiaggia nel sistema.
			\item[•] Per poter passare allo stato ATTIVO, la spiaggia deve possedere almeno i requisiti minimi operativi (Inventario, Servizi e Parcheggi definiti; almeno una Stagione e una Zona configurate, come definito dalle regole di business generali).
		\end{enumerate}
		
		\begin{enumerate}[leftmargin=*]
			\item[•] \textbf{Testing}: le classi che si occupano dei test sono \texttt{CreateBeachServiceTest} e \texttt{AddressServiceTest} nel package \texttt{service}, tutte le classi nei packages \texttt{domain.beach} e \texttt{domain.common}.
		\end{enumerate}
	}
		
	\newpage
	
		\usecasetable
	{UC-11}
	{Manage Beach}
	{Owner}
	{Permette all'Owner di aggiornare i dati generali, l'inventario, i servizi e i parcheggi, nel rigoroso rispetto dello stato di attività della spiaggia.}
	{L'Owner ha effettuato il login, possiede una spiaggia associata e l'account/spiaggia non sono soggetti a Ban (già verificato al login).}
	{Il database viene aggiornato con le nuove informazioni della spiaggia, dell'inventario, dei servizi e/o dei parcheggi.}
	{
		\begin{enumerate}[nosep, leftmargin=*]
			\item L'Owner accede alla dashboard di Gestione Spiaggia;
			\item Il sistema recupera i dati attuali e verifica lo stato della spiaggia (ATTIVO o INATTIVO);
			\item Il sistema "blocca" (sola lettura) o abilita i campi di modifica in base allo stato attuale;
			\item L'Owner naviga tra le sezioni e applica le modifiche desiderate (Anagrafica, Inventario, Servizi, Parcheggi);
			\item L'Owner clicca su "Salva Modifiche";
			\item Il sistema esegue i controlli di integrità (es. validazione dell'inventario contro le esigenze delle stagioni future e capienza parcheggi per le prenotazioni attive);
			\item Il sistema salva i dati aggiornati nel database e mostra un messaggio di successo.
		\end{enumerate}
	}
	{
		\begin{enumerate}[nosep, leftmargin=25px]
			\item[\textbf{3a.}] \textbf{Spiaggia in stato ATTIVO (Restrizioni):} Il sistema rende in sola lettura: Nome, Descrizione, Numero di Telefono, Indirizzo, Inventario, Servizi e Parcheggio. Non è permessa alcuna alterazione strutturale finché la spiaggia riceve prenotazioni.
			\item[\textbf{3b.}] \textbf{Spiaggia in stato INATTIVO (Permessi Estesi):} Il sistema sblocca la modifica dei dati base, Inventario, Servizi e Parcheggio.
			\item[\textbf{6a-e.}] \textbf{Conflitto Inventario:} Se la spiaggia è INATTIVA e l'Owner riduce la quantità di inventario (ombrelloni, tende), il sistema verifica le stagioni attuali e future. Se la nuova capacità è inferiore a quella utilizzata nelle Stagioni, il sistema blocca il salvataggio mostrando: "Inventario insufficiente: capacità ridotta sotto la soglia richiesta dalle stagioni programmate."
			\item[\textbf{6b-e.}] \textbf{Conflitto Parcheggi:} Se l'Owner riduce il numero di parcheggi, il sistema somma i posti auto richiesti da tutte le prenotazioni PENDING e CONFIRMED per ogni singola data (odierna e futura). Se il nuovo limite inserito è inferiore al totale già prenotato in una qualsiasi di queste date, il sistema blocca l'operazione mostrando: "Capacità parcheggio insufficiente per coprire le prenotazioni già attive o in attesa."
		\end{enumerate}
	}
	{
		\begin{enumerate}[nosep, leftmargin=*]
			\item[•] Dati anagrafici, Inventario, Servizi e Parcheggio sono modificabili esclusivamente in stato INATTIVO.
			\item[•] I controlli di capienza (Inventario e Parcheggi) devono garantire la retrocompatibilità sia con le Stagioni in corso e quelle in futuro che con le prenotazioni già accettate o in approvazione (CONFIRMED/PENDING).
		\end{enumerate}
		
		\begin{enumerate}[leftmargin=*]
			\item[•] \textbf{Testing}: le classi che si occupano dei test sono \texttt{ManageBeachServiceTest} nel package \texttt{service}; \texttt{BeachTest}, \texttt{BeachGeneralTest}, \texttt{BeachServicesTest}, \texttt{BeachInventoryTest} e \texttt{ParkingTest} nel package \texttt{domain}.
		\end{enumerate}
	}
		
	\newpage
	
	\usecasetable
	{UC-12}
	{Activate/Deactivate Beach}
	{Owner}
	{Permette all'Owner di modificare lo stato operativo della propria struttura, rendendola visibile (ATTIVA) o nascondendola (INATTIVA) per le nuove prenotazioni.}
	{L'Owner ha effettuato il login. Possiede una spiaggia nel sistema e questa non è soggetta a un Ban da parte dell'Admin.}
	{Lo stato della spiaggia viene aggiornato nel database. Se attivata, appare nelle ricerche dei Customer. Se disattivata, scompare dalle ricerche ma le prenotazioni attuali restano invariate.}
	{
		\begin{enumerate}[nosep, leftmargin=*]
			\item L'Owner accede alla dashboard principale della propria spiaggia;
			\item Il sistema mostra lo stato attuale della struttura tramite un apposito pulsante/interruttore (Es. "Attiva Spiaggia" o "Disattiva Spiaggia");
			\item L'Owner clicca sul pulsante per invertire lo stato attuale;
			\item Il sistema identifica il tipo di operazione richiesta (Attivazione o Disattivazione);
			\item Se l'operazione è di Attivazione, il sistema verifica la presenza dei requisiti minimi operativi (Inventario, Servizi, Parcheggi, Zone e Stagioni);
			\item Il sistema aggiorna lo stato nel database;
			\item Il sistema mostra un messaggio di conferma dell'avvenuto cambio di stato.
		\end{enumerate}
	}
	{
		\begin{enumerate}[nosep, leftmargin=25px]
			\item[\textbf{4a.}] \textbf{Flusso di Disattivazione:} Se l'Owner sceglie di disattivare la spiaggia, il sistema non richiede alcun controllo bloccante. Procede al punto 6, avvisando l'utente che la struttura non riceverà nuove prenotazioni e che quelle già registrate restano valide e gestibili tramite la sezione "Gestione Prenotazioni".
			\item[\textbf{5a-e.}] \textbf{Requisiti Mancanti per Attivazione:} Se l'Owner tenta di attivare la spiaggia ma il sistema rileva dati mancanti, l'operazione viene bloccata. Il sistema mostra un alert specifico: "Impossibile attivare la struttura. È necessario prima configurare:[Elenco dei dati mancanti, es. Inventario, almeno 1 Zona, almeno 1 Stagione futura/in corso].".
		\end{enumerate}
	}
	{
		\begin{enumerate}[nosep, leftmargin=*]
			\item[•] Per l'Attivazione sono richiesti obbligatoriamente: campi Inventario, campi Servizi, campi Parcheggio, almeno una Zona definita e almeno una Stagione configurata con data di fine $>$ data odierna.
			\item[•] La Disattivazione è un'operazione non distruttiva: non altera lo stato (PENDING/CONFIRMED) delle prenotazioni esistenti e permette all'Owner di continuare a operare su di esse.
			\item[•] Se l'intento è quello di chiudere definitivamente la spiaggia, vedi UC-13 (Close Beach).
		\end{enumerate}
		
		\begin{enumerate}[leftmargin=*]
			\item[•] \textbf{Testing}: le classi che si occupano dei test sono \texttt{ManageBeachServiceTest} nel package \texttt{service}; \texttt{BeachTest} nel package \texttt{domain}.
		\end{enumerate}
	}
		
	\newpage
	
	\usecasetable
	{UC-13}
	{Close Beach}
	{Owner}
	{Permette all'Owner di terminare definitivamente le operazioni della spiaggia, cancellando automaticamente tutte le prenotazioni future (PENDING/CONFIRMED) e notificando i clienti via email.}
	{L'Owner ha effettuato il login e la spiaggia è attualmente in stato ATTIVO.}
	{La spiaggia passa in stato CLOSED (Chiuso Definitivamente). Tutte le prenotazioni future risultano CANCELLED. Gli utenti interessati ricevono una notifica email.}
	{
		\begin{enumerate}[nosep, leftmargin=*]
			\item L'Owner accede alla dashboard e seleziona "Chiudi Attività";
			\item Il sistema mostra un avviso di sicurezza: "Attenzione, questa operazione chiuderà DEFINITIVAMENTE la spiaggia ed annullerà TUTTE le prenotazioni future (inviando una mail ai clienti). Questa azione NON È REVERSIBILE. Confermi?";
			\item L'Owner conferma l'operazione;
			\item Il sistema filtra tutte le prenotazioni (PENDING/CONFIRMED) con data $>$ data odierna;
			\item Il sistema imposta lo stato di tali prenotazioni a "CANCELLED";
			\item Il sistema invia una email automatica a ogni cliente coinvolto;
			\item Il sistema imposta la spiaggia in stato CLOSED.
		\end{enumerate}
	}
	{
		\begin{description}[nosep, leftmargin=25px]
			\item[\textbf{3a-e.}] \textbf{Conferma mancata:} Se l'Owner chiude la finestra di conferma, l'operazione viene annullata senza alcuna modifica ai dati.
		\end{description}
	}
	{
		\begin{enumerate}[nosep, leftmargin=*]
			\item[•] La chiusura è un'azione irreversibile riguardante le prenotazioni e la spiaggia. 
			\item[•] Vengono aggiornate nello stato CANCELLED solo le prenotazioni che devono ancora avvenire.
		\end{enumerate}
		
		\begin{enumerate}[leftmargin=*]
			\item[•] \textbf{Testing}: le classi che si occupano dei test sono \texttt{CloseBeachServiceTest} nel package \texttt{service}; \texttt{BeachTest} nel package \texttt{domain}.
		\end{enumerate}
	}
		
	\newpage
	
	\usecasetable
	{UC-14}
	{Manage Seasons}
	{Owner}
	{Permette all'Owner di definire i periodi operativi della spiaggia (date di inizio e fine) e di impostare contestualmente il listino prezzi base e le tariffe specifiche per le Zone incluse.}
	{L'Owner ha effettuato il login, l'account/spiaggia non sono soggetti a Ban. La spiaggia è in stato INATTIVO e non chiusa definitivamente.}
	{Il calendario delle stagioni, comprensivo di prezzi e tariffe zonali, viene aggiornato nel database.}
	{
		\begin{enumerate}[nosep, leftmargin=*]
			\item L'Owner accede alla dashboard e naviga nella sezione "Gestione Stagioni";
			\item Il sistema recupera e mostra l'elenco delle stagioni configurate;
			\item L'Owner seleziona "Nuova Stagione" (oppure "Modifica" / "Elimina" per una esistente);
			\item Il sistema applica i blocchi temporali e relazionali in base allo storico della stagione;
			\item L'Owner inserisce i parametri temporali (nome stagione, data inizio, data fine);
			\item L'Owner seleziona le Zone da includere e definisce i prezzi base al giorno (extra letini/sedie/sdraio, parcheggio, camerini) e le relative tariffe di Zona (prezzo ombrellone/tenda);
			\item L'Owner conferma l'operazione;
			\item Il sistema esegue le validazioni temporali (incluso il controllo anti-overlap) e salva i dati (date e prezzi) nel database.
		\end{enumerate}
	}
	{
		\begin{enumerate}[nosep, leftmargin=25px]
			\item[\textbf{3a.}] \textbf{Estensione Stagione (Date):} Se la stagione è in corso o futura, il sistema rende in sola lettura la Data di Inizio. L'Owner può solo posticipare la "Data di Fine".
			\item[\textbf{3b-e.}] \textbf{Prezzi Bloccati (Immutabilità Finanziaria):} Se l'Owner modifica una stagione, l'intera sezione relativa a prezzi base e tariffe di zona passa in modalità Sola Lettura.
			\item[\textbf{3c-e.}] \textbf{Eliminazione Bloccata:} Se si tenta di eliminare una stagione con prenotazioni associate, il sistema blocca l'azione mostrando l'avviso: "Impossibile eliminare: questa stagione contiene prenotazioni registrate.".
			\item[\textbf{8a-e.}] \textbf{Errore Sovrapposizione Date (Overlap):} Durante il salvataggio (creazione di una nuova stagione o posticipo della fine di una esistente), il sistema verifica che il nuovo intervallo di date non intersechi i periodi di altre stagioni. Se rileva un conflitto, blocca l'operazione mostrando l'errore: "Il periodo selezionato si sovrappone a una stagione esistente.", ed invita l'utente a modificare l'intervallo.
		\end{enumerate}
	}
	{
		\begin{enumerate}[nosep, leftmargin=*]
			\item[•] Le nuove stagioni possono essere create unicamente con date nel futuro rispetto alla data odierna.
			\item[•] Non è possibile accorciare la durata o posticipare l'inizio di una stagione già creata.
			\item[•] Prezzi e tariffe zonali diventano completamente immutabili dal secondo esatto in cui viene generata la prima prenotazione per quella specifica stagione.
		\end{enumerate}
		
		\begin{enumerate}[leftmargin=*]
			\item[•] \textbf{Mockup}: si veda Figura 19 - Gestione Stagioni e Prezzi
			
			\item[•] \textbf{Testing}: le classi che si occupano dei test sono \texttt{ManageBeachServiceTest} nel package \texttt{service}; \texttt{SeasonTest}, \texttt{ZoneTariffTest} e \texttt{PriceCalculatorTest} nel package \texttt{domain}.
		\end{enumerate}
	}
		
	\newpage
	
	\usecasetable
	{UC-15}
	{Manage Zones}
	{Owner}
	{Permette all'Owner di creare e configurare le diverse aree fisiche della spiaggia (es. Prima Fila, Area VIP) stabilendone il layout e i posti assegnati.}
	{L'Owner ha effettuato il login, l'account/spiaggia non sono soggetti a Ban. La spiaggia non è chiusa ed è attualmente in stato INATTIVO.}
	{Le zone e la loro configurazione spaziale vengono salvate nel database, pronte per essere incorporate all'interno delle Stagioni.}
	{
		\begin{enumerate}[nosep, leftmargin=*]
			\item L'Owner accede alla sezione "Gestione Zone";
			\item Il sistema mostra le zone attualmente presenti nel sistema;
			\item L'Owner seleziona "Nuova Zona", "Modifica" o "Elimina";
			\item Il sistema verifica se la zona selezionata possiede dipendenze (se è agganciata a una o più stagioni);
			\item L'Owner configura i dettagli (nome zona, posizionamento inventario/ombrelloni al suo interno) e clicca su "Salva";
			\item Il sistema registra la zona nel database.
		\end{enumerate}
	}
	{
		\begin{enumerate}[nosep, leftmargin=25px]
			\item[\textbf{4a-e.}] \textbf{Modifica/Eliminazione Bloccata (Zona In Uso):} Se l'Owner tenta di modificare o eliminare una zona che risulta già associata ad almeno una Stagione configurata (passata, presente o futura), il sistema inibisce l'azione rendendo i campi "Read-Only". Viene mostrato l'avviso: "Operazione negata. Questa zona è vincolata a una stagione e non può essere alterata."
		\end{enumerate}
	}
	{
		\begin{enumerate}[nosep, leftmargin=*]
			\item[•] Una volta creata, una zona è "libera" e completamente modificabile.
			\item[•] Nel momento esatto in cui la Zona viene aggiunta al set di una Stagione, questa diventa rigorosamente Immutabile in termini di layout, per preservare lo storico della disposizione della spiaggia.
		\end{enumerate}
		
		\begin{enumerate}[leftmargin=*]
			\item[•] \textbf{Mockup}: si veda Figura 20 - Gestione Zone e Layout
			
			\item[•] \textbf{Testing}: le classi che si occupano dei test sono \texttt{ManageBeachServiceTest} nel package \texttt{service}; \texttt{ZoneTest} e \texttt{SpotTest} nel package \texttt{domain}.
		\end{enumerate}
	}
		
	\newpage
	
	\usecasetable
	{UC-16}
	{Create Owner Account}
	{Admin}
	{Permette all'Admin di registrare manualmente un nuovo proprietario di spiaggia (Owner) nel sistema, assegnandogli le credenziali di primo accesso.}
	{L'Admin deve aver effettuato il login al proprio account.}
	{Viene creato un nuovo record utente di tipo "Owner" nel database. L'account viene contrassegnato con un flag che impone il cambio password al primo accesso.}
	{
		\begin{enumerate}[nosep, leftmargin=*]
			\item L'Admin riceve i dati del nuovo cliente (tramite telefono o email esterna);
			\item L'Admin accede alla propria dashboard e naviga nella sezione "Gestione Utenti";
			\item Clicca sul pulsante "Crea un nuovo utente Owner";
			\item Il sistema mostra il modulo di registrazione;
			\item L'Admin compila i campi obbligatori: Nome, Cognome, Username, Email e inserisce una Password temporanea;
			\item L'Admin clicca su "Salva";
			\item Il sistema verifica l'univocità dello Username e dell'Email;
			\item Il sistema salva il nuovo Owner nel database e mostra un messaggio di successo all'Admin.
		\end{enumerate}
	}
	{
		\begin{enumerate}[nosep, leftmargin=25px]
			\item[\textbf{7a-e.}] \textbf{Errore Dati Duplicati:} Se il sistema rileva che l'Email o lo Username inseriti appartengono già a un altro account (di qualsiasi ruolo) nel database, blocca il salvataggio. Evidenzia i campi incriminati e mostra l'errore: "Impossibile creare l'utente: Username o Email già registrati nel sistema."
		\end{enumerate}
	}
	{
		\begin{enumerate}[nosep, leftmargin=*]
			\item[•] Tutti i campi anagrafici e le credenziali sono obbligatori (non null).
			\item[•] Email e Username devono essere globalmente univoci all'interno di tutto il database utenti.
			\item[•] La password assegnata è temporanea (One-Time Password) e il nuovo Owner sarà forzato a cambiarla al primo accesso.
		\end{enumerate}
		
		\begin{enumerate}[leftmargin=*]
			\item[•] \textbf{Testing}: le classi che si occupano dei test sono \texttt{CreateUserServiceTest} e \texttt{CreateOwnerServiceTest} nel package \texttt{service}; \texttt{UserTest} e \texttt{OwnerTest} nel package \texttt{domain}.
		\end{enumerate}
	}
		
	\newpage
	
	\usecasetable
	{UC-17}
	{Create Admin Account}
	{Admin}
	{Permette ad un Admin di creare un nuovo account con privilegi di amministrazione (Admin), assegnandogli le credenziali di primo accesso.}
	{L'Admin creatore deve aver effettuato regolarmente il login al sistema.}
	{Viene creato un nuovo record utente di tipo "Admin" nel database, configurato con una One-Time Password per il primo accesso.}
	{
		\begin{enumerate}[nosep, leftmargin=*]
			\item L'Admin accede alla propria dashboard e naviga nella sezione "Gestione Utenti";
			\item Clicca sul pulsante "Crea nuovo account Admin";
			\item Il sistema mostra il modulo di registrazione dedicato;
			\item L'Admin compila tutti i campi obbligatori: Nome, Cognome, Username, Email e inserisce una Password temporanea;
			\item L'Admin clicca su "Salva";
			\item Il sistema verifica l'univocità dello Username e dell'Email nel database;
			\item Il sistema salva il nuovo account Admin e mostra un messaggio di conferma dell'avvenuta creazione.
		\end{enumerate}
	}
	{
		\begin{enumerate}[nosep, leftmargin=25px]
			\item[\textbf{6a-e.}] \textbf{Errore Dati Duplicati:} Se il sistema rileva che l'Email o lo Username inseriti sono già associati a un account esistente (di qualsiasi ruolo), blocca il salvataggio. Il sistema evidenzia i campi errati e mostra l'avviso: "Impossibile creare l'amministratore: Username o Email già presenti nel sistema."
		\end{enumerate}
	}
	{
		\begin{enumerate}[nosep, leftmargin=*]
			\item[•] Tutti i campi anagrafici e le credenziali sono obbligatori (non null).
			\item[•] Email e Username devono essere globalmente univoci all'interno dell'intero database utenti.
			\item[•] La password assegnata è temporanea (One-Time Password) e il nuovo Admin sarà forzato a cambiarla al primo accesso.
		\end{enumerate}
		
		\begin{enumerate}[leftmargin=*]
			\item[•] \textbf{Testing}: le classi che si occupano dei test sono \texttt{CreateUserServiceTest} e \texttt{CreateAdminServiceTest} nel package \texttt{service}; \texttt{UserTest} e \texttt{AdminTest} nel package \texttt{domain}.
		\end{enumerate}
	}
		
	\newpage
	
	\usecasetable
	{UC-18}
	{Manage Reports}
	{Admin}
	{Permette all'Admin di esaminare i Report in sospeso (creati da Customer o Owner) e decidere se approvarli o rifiutarli in base allo storico degli utenti coinvolti e alla descrizione del report.}
	{L'Admin ha effettuato il login al sistema. Esiste almeno un report in stato PENDING.}
	{Lo stato del report viene aggiornato in APPROVED o REJECTED e diventa immutabile.}
	{
		\begin{enumerate}[nosep, leftmargin=*]
			\item L'Admin accede alla propria dashboard e naviga nella sezione "Revisione Segnalazioni";
			\item Il sistema recupera e mostra l'elenco dei report in stato PENDING;
			\item L'Admin seleziona uno specifico report per aprirne il dettaglio;
			\item Il sistema mostra il testo della segnalazione, i dati del segnalante, i dati del segnalato e il numero totale di report pregressi (storico) a carico di entrambi gli utenti;
			\item L'Admin valuta la gravità della situazione e clicca su "Approva Segnalazione" oppure "Rifiuta Segnalazione";
			\item Il sistema salva il nuovo stato del report nel database;
			\item Se il report è stato approvato e l'Admin ritiene necessaria una sanzione, il flusso si sposta su UC-19 (Applica Ban).
		\end{enumerate}
	}
	{
		\begin{enumerate}[nosep, leftmargin=25px]
			\item[\textbf{2a.}] \textbf{Nessun report in sospeso:} Se non ci sono report PENDING, il sistema mostra il messaggio "Nessuna segnalazione da revisionare al momento."
			\item[\textbf{5a.}] \textbf{Rifiuto del Report:} Se l'Admin seleziona "Rifiuta Segnalazione", il processo termina al punto 6 senza alcuna conseguenza disciplinare per il segnalato.
		\end{enumerate}
	}
	{
		\begin{enumerate}[nosep, leftmargin=*]
			\item[•] Un report valutato (APPROVED o REJECTED) non può mai tornare in stato PENDING.
			\item[•] Il sistema deve obbligatoriamente mostrare lo storico dei report passati per fornire il contesto necessario all'Admin.
		\end{enumerate}
		
		\begin{enumerate}[leftmargin=*]
			\item[•] \textbf{Mockup}: si veda Figura 21 - Dashboard di Moderazione Admin 
			
			\item[•] \textbf{Testing}: le classi che si occupano dei test sono \texttt{ManageReportServiceTest} nel package \texttt{service}; \texttt{ReportTest} nel package \texttt{domain}.
		\end{enumerate}
	}
		
	\newpage
	
	\usecasetable
	{UC-19}
	{Apply Ban}
	{Admin}
	{Permette all'Admin di sanzionare in via definitiva un utente (Customer o Owner) o di interdirgli l'accesso a una specifica spiaggia, come conseguenza diretta di un report approvato.}
	{L'Admin sta eseguendo UC-18 (Manage Reports) ed ha appena approvato un report a carico dell'utente. L'utente non ha già subito la sanzione massima (App Ban).}
	{L'utente (o la spiaggia) viene bannato permanentemente. Le prenotazioni future coinvolte vengono cancellate in automatico.}
	{
		\begin{enumerate}[nosep, leftmargin=*]
			\item A seguito dell'approvazione del report (UC-18), l'Admin clicca sul pulsante "Applica Sanzione" relativo all'utente segnalato;
			\item Il sistema mostra le opzioni di Ban disponibili (App Ban o Beach Ban) in base al ruolo dell'utente;
			\item L'Admin seleziona il tipo di Ban desiderato e inserisce una motivazione formale;
			\item L'Admin clicca su "Conferma Ban Definitivo";
			\item Il sistema mostra un avviso critico (Warning): "Attenzione: l'azione è irreversibile. Confermi?";
			\item L'Admin conferma definitivamente;
			\item Il sistema aggiorna lo stato dell'utente o dei suoi permessi nel database;
			\item Il sistema procede alla cancellazione automatica delle prenotazioni future (come da regole di business) e notifica l'utente.
		\end{enumerate}
	}
	{
		\begin{enumerate}[nosep, leftmargin=25px]
			\item[\textbf{2a.}] \textbf{Opzioni limitate per Owner:} Se l'utente segnalato è un Owner, il sistema nasconde l'opzione "Beach Ban" e permette esclusivamente di applicare l'"App Ban".
			\item[\textbf{5a-e.}] \textbf{Annullamento operazione:} Se l'Admin annulla l'operazione al momento del warning, il Ban non viene applicato (ma il report rimane APPROVED).
		\end{enumerate}
	}
	{
		\begin{enumerate}[nosep, leftmargin=*]
			\item[•] I Ban sono definitivi. Non esiste alcuna procedura a sistema per la rimozione o la revoca di un Ban.
			\item[•] Conseguenze Beach Ban (Customer): Impedisce permanentemente la creazione di nuove prenotazioni nella spiaggia oggetto del report.
			\item[•] Conseguenze App Ban (Customer o Owner): Disattiva permanentemente l'accesso al sistema. Cancella automaticamente in stato CANCELLED tutte le prenotazioni future (PENDING o CONFIRMED). Se applicato a un Owner, disattiva d'ufficio e permanentemente la sua spiaggia.
		\end{enumerate}
		
		\begin{enumerate}[leftmargin=*]
		\item[•] \textbf{Mockup}: si veda Figura 21 - Dashboard di Moderazione Admin 
		
		\item[•] \textbf{Testing}: le classi che si occupano dei test sono \texttt{BanServiceTest} nel package \texttt{service}; \texttt{BanTest} nel package \texttt{domain}.
	\end{enumerate}
	}
\end{document}